\documentclass[12pt]{article}
\usepackage{geometry}
\geometry{letterpaper, margin=1in}
\usepackage{hyperref}
\usepackage[usenames,dvipsnames]{xcolor}
\usepackage{microtype}
\usepackage{enumitem}

\definecolor{accent}{RGB}{26, 60, 110}

\begin{document}

%----------HEADER----------
\begin{center}
  {\LARGE\textbf{\textcolor{accent}{SITONG LI}}} \\[4pt]
  \small Ann Arbor, MI
  \;\textbullet\;
  \href{mailto:ruc.sitongli@gmail.com}{ruc.sitongli@gmail.com}
  \;\textbullet\;
  (952) 201-4026
\end{center}

\vspace{1em}

\noindent May 22, 2026

\vspace{1em}

\noindent Dear USC Office of Institutional Research Hiring Committee,

\vspace{1em}

I am writing to apply for the Associate Data Scientist position in USC's Office of Institutional Research. I am currently a Data Scientist at the Ross School of Business, University of Michigan, where I build and maintain production ML pipelines and ETL workflows for large-scale administrative research. The work your office does --- translating complex institutional data into actionable insights that shape university strategy --- is a mission I find genuinely compelling, and it maps closely onto the data science infrastructure work I have been doing for the past several years.

My experience with data pipelines and ETL workflows spans several production deployments. At Michigan Ross, I built and own two end-to-end data systems: an automated entity resolution pipeline that integrates LLMs, string matching, and word embeddings to collect, transform, and validate records across heterogeneous administrative datasets, and a computer vision ML pipeline that performs structured attribute extraction at scale, with full data integrity checks throughout ingestion and output. For a research project at the University of Minnesota, I built a distributed web crawler to collect and parse 10,000+ full-text SEC filings from EDGAR, designed regex- and HTML-based segmentation pipelines, and then developed and optimized BERT-based ML classifiers across 500,000+ sentences --- all on datasets too large for in-memory processing, run on distributed infrastructure (Spark, Databricks, AWS). On a concurrent project, I built an automated ETL pipeline to extract structured information from 100,000+ regulatory documents across 1,000+ entities, deployed an agentic LLM with Google Search API integration for entity enrichment, and maintained comprehensive technical documentation throughout. In each case, the core discipline --- clean, validated, well-documented data pipelines that stakeholders can rely on --- is the same discipline your team needs.

Beyond pipeline engineering, I bring strong statistical and ML modeling skills and a consistent track record of communicating findings to non-technical audiences. I have built predictive models ranging from BERT-based text classifiers and attention-based pricing models to XGBoost feature selectors and word-embedding-based NLP models, and I have validated all of them through rigorous statistical testing (panel regressions, event study models, probit models). I hold a Master of Science in Finance and Business Analytics from the University of Minnesota's Carlson School (GPA 3.97), with coursework directly in Predictive Analytics, Causal Inference, and Big Data Analytics. I am proficient in Python and SQL, work daily in Git-based version-controlled environments, and have experience with Spark, Databricks, and AWS for large-scale data processing.

I am drawn to this role specifically because institutional research sits at the intersection of rigorous data practice and real organizational impact --- the kind of work where a well-built pipeline and a clearly communicated insight can directly inform how a university allocates resources and supports its community. I would welcome the opportunity to bring that orientation to USC's Office of Institutional Research.

\vspace{1em}

\noindent Sincerely,\\[0.5em]
\textbf{Sitong Li}

\end{document}
